\documentstyle[%


righttag%


,11pt%


,amssymb%


,russian%               remove in Eng.


,izvru%                 Set izven% in Eng.


]{amsart}





%





\newcommand{\tts}[1]{}





\theoremstyle{plain}


\theorembodyfont{\it}


\newtheorem{thmnonum}{\hskip\parindent Теорема}


\renewcommand{\thethmnonum}{}





%\hoffset=-15mm%                  For Eng.


\setcounter{page}{25}


\god{2005}


\nomer{1~(512)} %                        Remove in Eng.


%\nomer{Vol.\,49, No.\,1}%               Set in Eng.


%\str{\pageref{begin}--\pageref{end}}%  Set in Eng.


\udk{517.958}





\author{\it И.Б.\,Бадриев, О.А.\,Задворнов}


% NB:   ^^ remove in Eng.


\title{Исследование разрешимости осесимметричной задачи об


определении положения равновесия мягкой оболочки вращения}


\thanks{Работа выполнена при финансовой поддержке


Российского фонда фундаментальных исследований (коды проектов


\No\No\,03-01-00380, 04-01-00821) и Конкурсного центра фундаментального 


естествознания Министерства образования и науки Российской Федерации


(код проекта \No\,Е02-1.0-189).}





\date{}


%\pagestyle{myheadings}%         Set in Eng.


\pagestyle{plain} %              Remove in Eng.


\begin{document}


%\thispagestyle{plain}%          Set in Eng.


\maketitle


\label{begin}








\section{Введение}





Рассматривается осесимметричная задача об определении положения


равновесия мягкой оболочки вращения, образованной переплетением двух


семейств нитей, одно из


которых имеет циркулярное направление, а другое --- продольное.


Предполагается, что функция, определяющая в продольных нитях зависимость


модуля силы натяжения от степени удлинения, имеет линейный рост.


Ограничений на рост функции, определяющей в циркулярных нитях зависимость


модуля силы натяжения от степени удлинения, не накладывается.


Задача сформулирована математически в виде вариационного неравенства с


псевдомонотонным


оператором на замк\-нутом выпуклом множестве в гильбертовом пространстве.


Доказана теорема


существования. Ее доказательство основано на использовании теоремы о


неподвижной точке для многозначного оператора, определенного при помощи


вспомогательного вариационного неравенства.





Отметим, что ранее авторами изучались стационарные задачи теории мягких


оболочек (бесконечно длинных цилиндрических \cite{BS_91IV}


и сетчатых \cite{bz_Iz05_BZ_92IV}, \cite{bz_Iz05_Zad_2003_IzV}), в том


числе, и при


наличии препятствий. Для этих задач, сформулированных в виде операторных


уравнений, вариационных и квазивариационных неравенств, удалось установить


коэрцитивность в общепринятом смысле (см., напр.,


\cite{bz_Iz05_Lio_72},


\cite{bz_Iz05_ET_79}) входящих в уравнения и неравенства операторов, что


сделало возможным


применить для исследования их разрешимости общие результаты теории


псевдомонотонных операторов (см. \cite{bz_Iz05_Lio_72}).








\section{Постановка задачи}





Рассматривается осесимметричная задача о равновесии мягкой (т.\,е. не


воспринимающей


сжимающих усилий) сетчатой оболочки вращения, находящейся под воздействием


массовой и поверхностной нагрузок.


Оболочка образована переплетением двух семейств нитей,


одно из которых имеет радиальное направление, а другое --- продольное.


Края оболочки считаются закрепленными. Предполагается, что векторы


плотностей


поверхностных и массовых сил лежат в радиальной (проходящей через ось


симметрии)


плоскости, и перемещение точек оболочки происходит также в радиальном


направлении.


Поверхностная нагрузка предполагается следящей, т.\,е. направлена по нормали


к поверхности


оболочки. В недеформированном состоянии поверхность


оболочки представляет из себя цилиндр единичного радиуса и длиной $l$.


В качестве эйлеровой примем цилиндрическую систему координат $(\rho,


\varphi, z)$;


в силу осесимметричности задачи поверхность деформированной оболочки


описывается


координатами в продольном и радиальном направлениях $z=z(s)$,


$\rho =\rho(s)$,


$s$ --- лагранжева координата в продольном направлении. В недеформированном


состоянии $z_0(s)=s$, $\rho_0(s)=1$, $0 < s < l$.





Указанная задача в цилиндрической системе координат описывается следующей


системой дифференциальных уравнений \cite{bz_Iz05_R_Gul}:


\begin{align}


\label{shelBZB6}


\frac{d}{d s}


\bigg(\frac{T_{1}(\lambda_1)}{\lambda_1}\, \frac{d z}{ds} \bigg)


+ q \, \rho \frac{d \rho}{ds}+\wt{f_1} &= 0, \quad 0 < s < l;\\


%


\frac{d}{d s} \bigg(\frac{T_{1}(\lambda_1)}{\lambda_1}\frac{d \rho}{ds}


\bigg)- T_{2}(\lambda_2)-q \, \rho \frac{d z}{d s}+\wt{f_2} &= 0,


\quad 0 < s < l,


\label{shelBZB7}


\end{align}


где $\wt{f}$, $q$ --- известные функции, характеризующие массовые и


поверхностные силы,


$T_1$, $T_2$ --- функции, определяющие зависимость модуля силы натяжения в


нитях от относительных степеней удлинения $\lambda_1$, $\lambda_2$ в 


продольном и циркулярном


направлениях, $\lambda_1=((\rho')^2+(z')^2)^{\, 1/2}$, $\lambda_2 =\rho$.





Уравнения (\ref{shelBZB6}), (\ref{shelBZB7}) дополняются граничными


условиями


\begin{equation}


z(0)=0, \ \ z(l)=l, \ \ \rho(0)=1, \ \ \rho(l)=1.


\label{shelBZB8g}


\end{equation}


Кроме того, естественным для цилиндрической системы


координат является ограничение $\rho \geq 0$.





Вариационную задачу, соответствующую краевой задаче


(\ref{shelBZB6})--(\ref{shelBZB8g}), сформулируем


в перемещениях $u(s)=(u_1(s), u_2(s))$, $u_1=z(s)-s$,


$u_2(s) =\rho(s) -1$,


\begin{multline*}


\int_{0}^{l}\frac{T_{1}(\lambda_1)}{\lambda_1}


\big((1+u_1',u_2'), \eta' \big)ds+\int_{0}^{l}


q (1+ {u_2})\big(- u_{2}' 


\eta_{1}+({u}_{1}'+1) \eta_{2} \big)ds+\\


%


+ \int_{0}^{l}T_{2} (\lambda_2) \eta_2 ds


= \int_{0}^{l}(\wt{f}, \eta) ds


\quad \forall \eta \in C^{\infty}_{0}(0,l),


%% \label{shelBZB7_00}


\end{multline*}


где $\lambda_1=((1+u_1')^2+(u_2')^2)^{1/2}$, $\lambda_2=1+u_2$.





Относительно функций $T_1$ и $T_2$ считаем, что


\begin{gather}


T_i(\xi)=0, \ \ \xi \leq 1, \ \ i=1, 2


\ \ \text{(оболочка не воспринимает сжимающих усилий)},


\label{shelBZB8s} \\


%


T_i, \ \ i=1, 2, \ \ \text{непрерывныe, неубывающиe},


\label{shelBZB8v}


\end{gather}


$T_1$ имеет линейный рост на бесконечности, т.\,е. существуют


положительные $k_{0}$, $k_{1}$ такие, что


\begin{equation}


k_{0} (\xi-1) \leq T_1(\xi) \leq k_{1} \xi


\ \ \text{при} \ \ \xi \geq 1.


\label{shelBZB8t}


\end{equation}


Предполагаем также, что


\begin{equation}


q \in L_{\infty} (0,l).


\label{shelBZB8q}


\end{equation}


Введем пространство


$V=\big[\overset{\circ}{W}{}^{(1)}_2 (0, l)\big]^2$ с нормой


$$


\| u \|_V=\bigg({\int_{0}^{l} | u' |^2 ds}


\bigg)^{1/2},


$$


а также множество $K =\{u \in V : u_2+1 \geq 0$ $\forall s \in (0, l)\}$.


Очевидно, что множество $K$ является выпуклым и замкнутым.





Определим по функции $T_{1}$ вектор-функцию $g : R^{2} \to R^{2}$:


\begin{equation*}


g(y_1,y_2)=g(y) =


\begin{cases}


{T_{1}(| y |) | y |^{-1} y,} & {y \neq 0;} \\


{0,} & {y=0}.


\end{cases}


%% \label{Def_G}


\end{equation*}


Зададим операторы $A, B, D: V \to V$, а также элемент $f \in V$ формами


\begin{alignat*}{2}


(Au,\eta)&= \int_{0}^{l}


\big(g(1+u_1',u_2'), \eta' \big) ds, &\quad


(Du,\eta)&= \int_{0}^{l}T_{2} (1+u_2) \eta_2 ds, \\


%


(Bu,\eta)&= \int_{0}^{l}q (1+ {u_2})


\big(- {u}_{2}' \eta_{1}+({u}_{1}'+1) \eta_{2} \big)ds,


&\quad (f,\eta)&=\int_{0}^{l}(\wt f,\eta) ds.


\end{alignat*}


Корректность определения этих операторов следует из непрерывности


$T_2$, вложения $\overset{\circ}{W}{}^{(1)}_2 (0, l)$ в $C(0,l)$


и условий (\ref{shelBZB8t}), (\ref{shelBZB8q}).





Под обобщенным решением осесимметричной задачи об определении положения


равновесия мягкой оболочки вращения, закрепленной по краям, находящейся


под воздействием массовой и поверхностной нагрузок, будем понимать функцию


$u \in K$, удовлетворяющую вариационному неравенству


\begin{equation}


((A+D)u+B u, \eta-u)\geq (f, \eta-u) \quad \forall \, \eta \in K.


\label{shelBZB8}


\end{equation}








\section{Свойства операторов и существование решения}





Для исследования разрешимости задачи (\ref{shelBZB8}) установим некоторые


свойства операторов\linebreak $A$, $B$, $D$.





\begin{lem} \label{BZ_Izv05_L1}


Пусть выполнены условия $(\ref{shelBZB8s})$--$(\ref{shelBZB8t})$. Тогда


оператор $A$ является монотонным и ограниченным и, кроме того,


\begin{equation}


(Au,u) \geq


k_{0} \| u \|_V^2-(3 k_{0}+k_{1}) \sqrt {l} \, \| u \|_V -


k_{0} l \quad \forall \, u \in V.


\label{shelBZB8Ak}


\end{equation}


\end{lem}





\begin{pf}


Из определения функции $g$ в силу монотонного возрастания $T_1$ для любых


векторов $y, z \in R^{\,2}$ имеем


\begin{gather*}


(g(y)-g(z), y-z ) =


\bigg(\frac{T_1(|y|)}{|y|} y -


\frac{T_1(|z| )}{|z|} z, y-z \bigg) =\\


%


= T_1(|y| ) |y| -\bigg[ \frac{T_1(|y| )}{|y|} +


\frac{T_1(|z| )}{|z|} \bigg] (y, z)


+ T_1(|z| ) |z| \geq


\big[ T_1(|y| )-T_1(|z| )\big] (|y|-|z| ) \geq 0.


\end{gather*}


Отсюда получаем


\begin{gather*}


(Au-Av,u-v)= \int_{0}^{l}


\big(g(1+u_1',u_2')-g(1+v_1',v_2'), u'-v' \big) ds =\\


%


= \int_{0}^{l}


\big(g(1+u_1',u_2')-g(1+v_1',v_2'),


 ((1+u_1')-(1+v_1'),u_2' -u_2' ) \big) ds \geq 0


\ \; \forall u, v \in V,


\end{gather*}


т.\,е. $A$ --- монотонный оператор.





Далее, из (\ref{shelBZB8s}), (\ref{shelBZB8t}) вытекает


$k_{0} (\xi-1) \leq T_1(\xi) \leq k_{1} \xi$


$\forall \xi \geq 0$,


следовательно,


\begin{multline}


(g(y), z )=(g(y), y )+(g(y), z-y) \geq


T_{1}(| y | ) | y |-T_{1}(| y | ) | z-y | \geq \\


%


\geq k_{0} (| y |-1) | y |-k_{1} | y | | z-y | =


k_{0} | y |^2-(k_{0}+k_{1} | z-y | ) | y |.


\label{shelBZB8tt}


\end{multline}





Пусть $u$ --- произвольный элемент из $V$. Положим $y=(1+u_1',u_2')$,


$z=u'=(u_1',u_2')$. При этом


$| z-y |= 1$, а $|y| \leq |u' |+1$.


Подставляя указанные векторы $y$, $z$ в неравенство (\ref{shelBZB8tt}),


получаем


$$


\big(g(1+u_1',u_2'), u' \big) \geq


k_{0} \big((1+u_1')^2+(u_2' )^2 \big) -


(k_{0}+k_{1} ) |u' |^2 \geq


k_{0} |u' |^2-(3 k_{0}+k_{1} ) |u'|-k_{0},


$$


откуда вытекает 


$$


(Au,u)= \int_{0}^{l}\big(g(1+u_1',u_2'), u' \big) ds


\geq k_{0} \| u \|_V^2-(3 k_{0}+k_{1} ) 


\sqrt {l} \| u \|_V-k_{0} l\ \ \forall u \in V.


$$





Ограниченность оператора $A$ непосредственно следует из условия


(\ref{shelBZB8t}).


\end{pf}





\begin{lem} \label{BZ_Izv05_L2}


Пусть выполнены условия $(\ref{shelBZB8s})$, $(\ref{shelBZB8v})$. Тогда


оператор $D$ является монотонным, ограниченным, и, кроме того,


\begin{equation}


(Du,u) \geq 0 \quad \forall u \in V.


\label{shelBZB8Dk}


\end{equation}


\end{lem}





{\bf Доказательство.}


B силу монотонного возрастания $T_2$ получаем монотонность оператора $D$:


\begin{gather*}


(Du-Dv,u-v)=\int_{0}^{l}


\big[ T_{2}(1+u_2)-T_{2}(1+ v_2)\big](u_2-v_2) ds =\\


%


= \int_{0}^{l}\big[ T_{2}(1+u_2)-T_{2}(1+ v_2)\big] 


((1+u_2)-(1+v_2) ) ds\geq 0


\quad \forall u, v \in V.


\end{gather*}





Ограниченность оператора $D$ следует непосредственно из непрерывности


$T_2$ и вложения $\overset{\circ}{W}{}^{(1)}_2 (0, l)$ в $C(0,l)$.





Установим справедливость (\ref{shelBZB8Dk}). Пусть


$u$ --- произвольный элемент из $V$. Положим


$\Omega^+=\{s \in (0, l): u_2(s) \geq 0 \}$.


Вне множества $\Omega^+$ выполнено неравенство $1+u_2 \leq 1$ и в силу


условия (\ref{shelBZB8s}) вне $\Omega^+$ имеем $T_2(1+u_2)=0$,


следовательно,


$$


(Du,u)=\int_{0}^{l}T_{2}(1+u_2) u_2 ds =


\int_{\Omega^+}T_{2}(1+u_2) u_2 ds \geq 0


\quad \forall u \in V.\quad\square


$$








\begin{lem} \label{BZ_Izv05_L3}


Пусть выполнено условие $(\ref{shelBZB8t})$. Тогда, если


$u_n \rightharpoonup{u}$, $\eta_n \rightharpoonup \eta$ в $V$


при $n \to+\infty$, то


\begin{equation}


\lim_{n \to+\infty} (Bu_n,\eta_n)=(Bu,\eta),


\label{BZ_Izv05_L3_B_PsMon}


\end{equation}


оператор $B$ ограничен и является псевдомонотонным.


\end{lem}





\begin{pf}


Пусть $u_n \rightharpoonup{u}$, $\eta_n \rightharpoonup \eta$ в


$V$ при $n \to+\infty$. Тогда


\begin{gather*}


\big| {(B u_n, \eta_n )-(B u, \eta )} \big|=


\bigg| {\int_{0}^{l}


(1+ {u_{2n}}) q \big[-{u}_{2n}' 


\eta_{1n}+({u}_{1n}'+1) \eta_{2n}\big] ds}- \\


%


- {\int_{0}^{l}


(1+ u_{2}) q [-{u}_{2}' 


\eta_{1}+({u}_{1}'+1) \eta_{2}] ds} \bigg| \leq 


\bigg| {\int_{0}^{l}q {u}_{2n}' 


[(1+ {u_{2n}}) \eta_{1n} -(1+ {u_{2}}) \eta_{1}] ds} \bigg| + \\


%


+ \bigg| {\int_{0}^{l}q ({u}_{1n}'+1) 


[(1+ {u_{2n}}) \eta_{2n}-(1+ {u_{2}}) \eta_{2}] ds} \bigg| +\\


%


+ \bigg| {\int_{0}^{l}q(1+ {u_{2}}) \eta_{1} 


[ {u}_{2n}'-{u}_{2}' ] ds} \bigg| +


\bigg| {\int_{0}^{l}q (1+ {u_{2}}) \eta_{2} 


[ ({u}_{1n}'+1)-({u}_{1}'+1)] ds} \bigg| = 


\sum_{j=1}^{4} I_{j n}.


\end{gather*}





В силу компактного вложения $\overset{\circ}{W}{}^{(1)}_2(0, l)$


в $C(0,l)$ и $L_2(0,l)$ последовательности


$\{{u_{in}}\}_{n=1}^{\infty}$ и


$\{{\eta_{in}}\}_{n=1}^{\infty}$, $i=1, 2$, сходятся сильно в


$C(0,l)$ и $L_2(0,l)$ к $u_i$ и $\eta_i$, $i=1, 2$, соответственно


при $n \to+\infty$.





Поскольку $q \in L_{\infty}(0, l)$, а последовательности


$\{{u_{in}}\}_{n=1}^{\infty}$,


$i=1, 2$, ограничены в $\overset{\circ}{W}{}^{(1)}_2 (0, l)$


(в силу их слабой сходимости), то


$$


I_{in} \leq c_i \int_{0}^{l}|(1+ {u_{2n}}) \eta_{in} -


(1+ {u_{2}}) \eta_{i} | ds \to 0 \ \ \text{при}\ \ n \to+\infty,\ \ i=1, 2.


$$





Так как $q(1+ {u_{2}}) \eta_{i}$, $i=1, 2$, порождают


линейные непрерывные функционалы на $\overset{\circ}{W}{}^{(1)}_2 (0,l)$, то


$I_{3n} \to 0$, $I_{4n} \to 0$ при $n \to+\infty$,


тем самым справедливость (\ref{BZ_Izv05_L3_B_PsMon}) установлена.





Доказательство ограниченности $B$ устанавливается на основании


вложений $\overset{\circ}{W}{}^{(1)}_2 (0, l)$


в $C(0,l)$ и $L_2(0,l)$ и условия (\ref{shelBZB8q}).





Наконец, в силу (\ref{BZ_Izv05_L3_B_PsMon}) из того,


что $u_n \rightharpoonup u$ в $V$ при $n \to+\infty$ и


$\limsup\limits_{n \rightarrow+\infty}(Bu_n,u_n-u) \leq 0$,


следует неравенство


$\liminf\limits_{n \rightarrow+\infty}(Bu_n,u_n-\eta) \geq (Bu,u-\eta)$,


т.\,е $B$ --- псевдомонотонный оператор.


\end{pf}





Определим теперь многозначный оператор $G : V \rightarrow 2^V$, где $2^V$


--- множество всех подмножеств пространства $V$, сопоставляя каждому


$w \in V$ множество $G(w)$ всех решений вариационной задачи


\begin{equation}


u \in K : ((A+D)u, \eta-u)\geq (f- B w, \eta-u) \quad \forall 


\eta \in K.


\label{shelBZB8_01}


\end{equation}





\begin{lem} \label{BZ_Izv05_L4}


Пусть выполнены условия $(\ref{shelBZB8s})$--$(\ref{shelBZB8q})$. Тогда


при каждом $w \in V$ множество $G(w)$ не пусто, ограничено, выпукло,


замкнуто и, следовательно, отображение $G$ является строго многозначным


с выпуклыми, компактными в слабой топологии пространства $V$, значениями.


\end{lem}





\begin{pf}


Как уже отмечалось выше, множество $K$ является выпуклым и замкнутым.


Из лемм \ref{BZ_Izv05_L1}, \ref{BZ_Izv05_L2} следует, что оператор $A+D$


является псевдомонотонным, из неравенств (\ref{shelBZB8Ak}),


(\ref{shelBZB8Dk}) вытекает


\begin{equation}


((A+D)u,u) \geq c_3 \|u\|^2-c_4 \|u\|-c_5


\quad \forall u \in V,


\label{shelBZB8A+Dk}


\end{equation}


т.\,е. оператор $A+D$ является коэрцитивным, поэтому задача


(\ref{shelBZB8_01}) имеет по крайней мере одно решение


(см., напр., \cite{bz_Iz05_Lio_72}, с.\,261) для


любого $w \in V$, а значит, при каждом $w \in V$ множество $G(w)$ не


пусто.





Далее, в силу коэрцитивности $A+D$ из (\ref{shelBZB8_01}) при $\eta=0$


получаем, что при каждом $w \in V$ множество $G(w)$ ограничено.





Следуя \cite{bz_Iz05_KL_73_IV},


нетрудно проверить, что оператор $A+D$ является непрерывным.


Кроме того, в силу лемм \ref{BZ_Izv05_L1} и \ref{BZ_Izv05_L2} оператор


$A+D$ монотонен. Поэтому множество решений задачи (\ref{shelBZB8_01})


при каждом $w \in V$


выпукло и замкнуто (см., напр., \cite{bz_Iz05_Lio_72}, с.\,263).





Из рефлексивности пространства $V$ следует


вторая часть утверждения леммы.


\end{pf}





\begin{lem} \label{BZ_Izv05_L5}


Пусть выполнены условия $(\ref{shelBZB8s})$--$(\ref{shelBZB8q})$. Тогда


график $\mathrm{graph}\,(G)=\{(w,u) \in V \times V : w \in V$, $u \in


G(w)\}$ отображения $G$ секвенциально слабо замкнут.


\end{lem}





\begin{pf}


Пусть $w_n \rightharpoonup w$ в $V$ при $n \to+\infty$,


$u_n \in G w_n$, $u_n \rightharpoonup u$ в $V$ при $n \to+\infty$, т.\,е.


\begin{equation}


u_n \in K : ((A+D) u_n, \eta-u_n)\geq (f-B w_n, \eta-u_n)


\quad \forall \eta \in K.


\label{shelBZB8_02}


\end{equation}


Докажем, что $u \in G w$. Поскольку $u_n \in K$, как решение задачи


(\ref{shelBZB8_02}), а множество $K$ является выпуклым, замкнутым,


и, следовательно, слабо замкнутым, то $u \in K$.


В силу леммы~\ref{BZ_Izv05_L3} \ $\lim\limits_{n \to+\infty}


(f-B w_n, u_n-\eta)=(f-B w, u -\eta)$


для любого $\eta \in K$. С учетом этого соотношения из (\ref{shelBZB8_02})


с $\eta=u$ имеем


$$


\limsup_{n \to+\infty} ((A+D) u_n, u_n-u) \leq


\lim_{n \to+\infty} (f-B w_n, u_n-u)=0,


$$


поэтому в силу псевдомонотонности оператора $A+D$ из (\ref{shelBZB8_02})


получаем


$$


(f-B w, u-\eta)=\lim_{n \to+\infty} (f-B w_n, u_n -\eta)\geq


\liminf_{n \to+\infty} ((A+D) u_n, u_n-\eta) \geq((A+D) u, u-\eta),


$$


откуда вытекает, что $u$ --- решение вариационного неравенства


(\ref{shelBZB8_01}), т.\,е. $u \in G w$.


\end{pf}





\begin{thmnonum} \label{BZ_Izv05_T1}


Пусть выполнены условия $(\ref{shelBZB8s})$--$(\ref{shelBZB8q})$.


Тогда для любого $c{>}0$ существует такое $\delta_c{>}0$,


что задача $(\ref{shelBZB8})$ имеет по крайней мере одно


решение при условиях


$\| f \| \leq c$, $q^*=\| q \|_{L_{\infty}(0,l)} \leq \delta_c$.


\end{thmnonum}





%%%%%% remove extra { } ^^^^^^








\begin{pf}


Для произвольного $R > 0$ введем выпуклое, замкнутое,


ограниченное множество $K_R=K \bigcap S_R$, где


$S_R= \{u \in V : \|u\|\leq R \}$.


Докажем, что для произвольного $c>0$ существуют такие $R>0$ и $\delta_c>0$,


при которых оператор $G$ отображает множество $K_R$ в себя.





Пусть $w$ --- произвольный элемент из $K_R$. Очевидно, что $0 \in K_R$.


Тогда, подставляя в (\ref{shelBZB8_01}) $\eta=0$, получим


$(Au+Du,-u) \geq (f, -u)-(Bw,-u)$,


откуда


$$


(Au+Du,u) \leq (f, u)-(Bw,u) \leq \| f \|\, \|u\|+ q^* c_6 \|w\|^2 


\|u\|+q^* c_7 \|w\|\, \|u\|.


$$


Отсюда в силу (\ref{shelBZB8A+Dk})


$$


c_3 \|u\|^2-b \|u\|-c_5 \leq 0, \quad b=c_4 +


q^* c_6 \|w\|^2+q^* c_7 \|w\|+\| f \|,


$$


следовательно, $ \|u\| \leq \rho (q^*, \| f \|, \|w\|)$, где


$$


\rho (q^*, \| f \|, \|w\|)=\frac{b +


\sqrt{b^2+ 4 c_3 c_5}}{2 c_3}.


$$





Далее, для $w \in K_R$ при $\| f \| \leq c$ имеем


\begin{multline*}


\|u\| \leq \rho^*(q^*)=


\max_{\| f \| \leq c,\; \| w \| \leq R} \rho(q^*,\| f \|,\|w\|)=\\


%


= \frac{c_4+q^* c_6 R^2+q^* c_7 R+c +


\sqrt{(c_4+q^* c_6 R^2+q^* c_7 R+c)^2+ 4 c_3 c_5}}{2 c_3}.


\end{multline*}





Положим


$$


R= \frac{c_4+ c+\sqrt{(c_4+c)^2+ 4 c_3 c_5}}{c_3}.


$$


Тогда при


$$


0 \leq q^* \leq \frac {c_4+c}{c_6 R^2+c_7 R}=\delta_c


$$


имеем $q^* c_6 R^2+q^* c_7 R \leq c_4+ c$, следовательно,


$$


\|u\| \leq \rho^*(q^*) \leq


\frac{2 (c_4+ c)+\sqrt{4 (c_4+c)^2+ 4 c_3 c_5}}{2 c_3}\leq R,


$$


а значит, $u \in K_R$.





Таким образом, оператор $G$ при некотором


$R$ и достаточно малом $q^*$ переводит множество $K_R$ в себя.


Поскольку $V$ рефлексивно и сепарабельно, то множество $K_R$,


рассматриваемое с индуцированной в нем слабой топологией пространства $V$,


является метрическим компактом. Из леммы \ref{BZ_Izv05_L5} следует,


что многозначное отображение $G : K_R \rightarrow K_R$ обладает


замкнутым графиком, и, следовательно, полунепрерывно сверху


(см. \cite{bz_Iz05_OE_88}, с.\,117). Поэтому, в силу


теоремы Какутани (\cite{bz_Iz05_OE_88}, с.\,336) $G$ имеет по крайней


мере одну неподвижную точку $u^*$ из $K_R$: $G u^*=u^*$. Подставляя в


(\ref{shelBZB8_01}) \ $u=w=u^*$, получим, что $u^*$ --- решение


вариационного неравенства (\ref{shelBZB8}).


\end{pf}








\begin{thebibliography} {9}





\bibitem{BS_91IV}


Бадриев И.Б., Шагидуллин Р.Р.


{\em Исследование одномерных уравнений статического


состояния мягкой оболочки и алгоритма их решения} // Изв.


вузов. Математика. -- 1992. -- \No\,1. -- С.\,7--17.





\bibitem{bz_Iz05_BZ_92IV}


Бадриев И.Б., Задворнов О.А.


{\em Исследование разрешимости стационарных задач для


сетчатых оболочек} // Изв. вузов. Математика. -- 1992.


-- \No\,11. -- С.\,3--7.





\bibitem{bz_Iz05_Zad_2003_IzV}


Задворнов О.А.


{\em Постановка и исследование стационарной задачи о контакте


мягкой оболочки с препятствием} // Изв. вузов. Математика.  -- 


2003. -- \No\,1. -- С.\,45--52.





\bibitem{bz_Iz05_Lio_72}


Лионс Ж.-Л.


{\em Некоторые методы решения нелинейных краевых задач}.  -- 


М.: Мир, 1972. -- 588~с.





\bibitem{bz_Iz05_ET_79}


Экланд И., Темам Р. {\em Выпуклый анализ и вариационные проблемы}.  -- 


М.: Мир, 1979. -- 400~c.





\bibitem{bz_Iz05_R_Gul}


Ридель В.В., Гулин Б.В.


{\em Динамика мягких оболочек}. --  М.: Наука, 1990. -- 206~с.





\bibitem{bz_Iz05_KL_73_IV}


Ляшко А.Д., Карчевский М.М.


{\em О решении некоторых нелинейных задач теории фильтрации}


// Изв. вузов. Математика. -- 1975. -- \No\,6. -- С.\,73--81.





\bibitem{bz_Iz05_OE_88}


Обэн Ж.-П., Экланд И.


{\em Прикладной нелинейный анализ}. --  М.: Мир, 1988. --  516~с.





\end{thebibliography}





\vskip 2cm





\post{\it Казанский государственный}{\it Поступила}


\post{\it университет}{31.08.2004}





\label{end}


\end{document}


